\section{Actual quadruple cycles and their local unfolding}
\label{sec:actual-bifurcations}

We use the fixed full-rank observation inverse with $\xi=\xi_0$,
$A_1(0)=B>0$ and $c_0=-C<0$. The extremal quartic continues to an
actual quadruple cycle at positive scales. To obtain its local
unfolding, we first localize all roots on the common interval, including
across the moving events, and then apply smooth preparation.

\begin{lemma}[Main-scale values across the events]
\label{mainloc:cross-event}
Let $A_1(0)=B>0$, $c(0)=-C<0$, and set
\begin{equation}\label{mainloc:weights}
 a=\frac{m^4A}{C^3},\qquad b=-\frac{m^2S}{C},\qquad
 \delta=\frac{2m^3T}{3BC^2},\qquad q=\frac{m^2\zeta}{C^2},
 \qquad m>0.
\end{equation}
For every compact control set and every bounded $\zeta$ interval,
uniformly on their product,
\begin{equation}\label{mainloc:full-limit}
 \frac{C^3}{m^4}D\left(\frac{m^2\zeta}{C^2}\right)
 =A-S\zeta-\zeta^2+\zeta_+^{3/2}
                 +T\sqrt{\zeta_+}+O(\sqrt m).
\end{equation}
This is a $C^0$ estimate, including both actual events.
\end{lemma}

\begin{proof}
Hold $a,b,m,\xi_0$ fixed and integrate in the separation control
$t$ from $0$ to $\delta$. Put $y=\Gamma_1(q,t)-t$. The passage
jet and its parameter derivatives give
\[
 \Gamma_{\nu,q}(x)=1+\tfrac32A_\nu\sqrt{x_+}+O(x_+),\qquad
 \Gamma_{\nu,t}(x)=O(x_+^{3/2}),\qquad
 |y-q|\le K(q_+^{3/2}+|t|).
\]
These estimates include zero inputs. Along the integration segment,
$A_2=-B+O(m)$ and $|t|=O(m^3)$. The square-root inequality
$|\sqrt{y_+}-\sqrt{q_+}|\le\sqrt{|y-q|}$ and the total derivative
\eqref{eq:quartic-total-detuning} therefore give
\begin{align*}
 \left|Q_t-\tfrac32B\sqrt{q_+}\right|
 \le K\{&q_+^{3/4}+\sqrt{|t|}+m\sqrt{q_+}+q_++|t|\\
        &+q_+^{3/2}+(q_++|t|)^{3/2}\}=O(m^{3/2}).
\end{align*}
This bound uses $|t|$ and remains valid for $q<0$, so both signs
of separation are included. Integration and the aligned identity
\eqref{eq:passage:aligned-jet} yield
\[
 Q-q=mq_+^{3/2}+\tfrac32B\delta\sqrt{q_+}+O(m^{9/2}),
 \qquad Q-q=O(m^4).
\]
In \eqref{eq:quartic-closing-division}, $c=-C+O(m)$ because the
fixed full-rank observation inverse has $\xi=\xi_0$. Replacing $Q$ by $q$
in $bQ+cQ^2$ changes these terms by $O(m^6)$, while
$(c+C)q^2=O(m^5)$ and $Q^3S=O(m^6)$. Thus
\[
 D=a+bq-Cq^2+mq_+^{3/2}
                    +\tfrac32B\delta\sqrt{q_+}+O(m^{9/2}).
\]
Substitution of \eqref{mainloc:weights} proves the assertion.
\end{proof}

With \eqref{mainloc:weights}, set $\Theta=(A,S,T)$ and
\begin{equation}\label{quartic:centre}
 \Theta_0=\left(-\frac1{256},\frac38,\frac1{16}\right),\qquad
 x_0=\frac14.
\end{equation}
To count cycles on both sides of this bifurcation, we first locate
all zeros without a saturation assumption.

\begin{lemma}[All roots near the quadruple configuration]
\label{mainloc:all-roots}
There are a fixed control neighbourhood $U$ of $\Theta_0$, $m_0>0$
and three disjoint compact state intervals $J_{N,-},J_{N,+},J_Q$
about, respectively,
\[
 \zeta_{N,\pm}=\frac{-3\pm2\sqrt2}{16}<0,\qquad
 \zeta_Q=\frac1{16}>0,
\]
such that, for every $\Theta\in U$ and $0<m<m_0$, all zeros in
$I$ lie in $(m^2/C^2)(J_{N,-}\cup J_{N,+}\cup J_Q)$.
Each negative interval contains exactly one simple zero. The positive
interval is both-active and contains all remaining zeros.
The three intervals can be chosen arbitrarily small about their
respective centres, by reducing $U$ and $m_0$.
\end{lemma}

\begin{proof}
Bounded controls satisfy the hypotheses of Lemma~\ref{mainloc:tail}.
Thus all zeros have $|\zeta|\le L$ for a fixed $L$.
At $\Theta_0$, the full-state limit in
\eqref{mainloc:full-limit} is
\[
 -\left(\zeta-\zeta_{N,-}\right)
       \left(\zeta-\zeta_{N,+}\right)\quad(\zeta<0),\qquad
 -\left(\sqrt\zeta-\tfrac14\right)^4\quad(\zeta\ge0).
\]
Its only root positions are the three stated centres, and its value
at zero is $-1/256$. Choose the three small intervals. On their
compact complement in $[-L,L]$, the absolute value of the limiting
function has a positive minimum. First reducing $U$, then $m_0$,
preserves that gap by Lemma~\ref{mainloc:cross-event}. This excludes
all zeros outside the intervals, including the event layer.

Take $U$ so that $T>0$. For $q<0$ both passages are inactive, so
the closing expansion gives on either negative interval
\[
 \frac{C^3}{m^4}D(m^2\zeta/C^2)
 =A-S\zeta-\zeta^2+O_{C^1}(m).
\]
The limiting derivatives at $\zeta_{N,-},\zeta_{N,+}$ are
$\sqrt2/4,-\sqrt2/4$. Their strict signs and the opposite endpoint
values persist uniformly. The intermediate value theorem and strict
monotonicity give one simple root in each interval. In particular,
their scaled locations are
\[
 \frac{-S\pm\sqrt{S^2+4A}}2+O(m).
\]
Finally the true second event is
$e=\Gamma_1^{-1}(\delta)=O(m^3)$, whereas the positive interval
has $q\ge k m^2$ for a fixed $k>0$. Hence that entire interval
lies beyond both events.
\end{proof}

\begin{theorem}[Actual quadruple-cycle bifurcation]
\label{thm:actual-quartic}
After a common reduction of the fixed control neighbourhood $U$ and
$m_0$ in Lemma~\ref{mainloc:all-roots}, the following hold in $U$
for every $0<m<m_0$. There is a unique local smooth
branch $(\Theta_*(m),x_*(m))$ through $(\Theta_0,x_0)$ at $m=0$
on which the actual displacement has a zero of ordinary multiplicity
four. Its physical coordinates satisfy
\begin{equation}\label{quartic:physical-branch}
 \begin{aligned}
 a_*(m)&=-\frac{m^4}{256C^3}+O(m^5),&
 b_*(m)&=-\frac{3m^2}{8C}+O(m^3),\\
 \delta_*(m)&=\frac{m^3}{24BC^2}+O(m^4),&
 q_*(m)&=\frac{m^2}{16C^2}+O(m^3).
 \end{aligned}
\end{equation}
There are smooth actual control coordinates $(u,v,w)$ centred on
this branch, and a smooth state coordinate $y$, in which the local
zero equation is, up to a nonvanishing smooth factor,
\begin{equation}\label{quartic:preparation}
 p(y)=y^4+uy^2+vy+w=0.
\end{equation}
In the normalized variables $(m,x,A,S,T)$, the changes and the
nonzero factor extend jointly smoothly through $m=0$. Away from the multiple-root set, the total
number of cycles in $I$ is two, four or six; each count occurs on
a nonempty open set of actual controls. At the quadruple point
there are exactly three distinct cycles: the quadruple cycle and
two simple negative-state companions.
\end{theorem}

\begin{proof}
Lemma~\ref{lem:weighted-quartic}, after constant rescaling, gives the
jointly smooth auxiliary function
\begin{equation}\label{quartic:joint-quotient}
 F(m,x;\Theta)=\frac{C^3}{m^4}D(m^2x^2/C^2)
 =A-Sx^2+x^3-x^4+Tx+mE_4(m,x,\Theta)
\end{equation}
near $(0,x_0,\Theta_0)$, with all prescribed finite mixed derivatives
bounded. Physical equality holds for $m>0$.
For $H=(F,F_x,F_{xx},F_{xxx})$, differentiation at that point,
in columns $(A,S,T,x)$, gives
\[
 D_{(A,S,T,x)}H=
 \begin{pmatrix}
 1&-1/16&1/4&0\\
 0&-1/2&1&0\\
 0&-2&0&0\\
 0&0&0&-24
 \end{pmatrix},\qquad \det D_{(A,S,T,x)}H=-48.
\]
The implicit function theorem gives the unique local branch
$\Theta_* =\Theta_0+O(m)$, $x_*=x_0+O(m)$, with
$F_{xxxx}=-24+O(m)\ne0$. Substituting the weights gives
\eqref{quartic:physical-branch}. Lemma~\ref{mainloc:all-roots}
places this zero in the both-active window. For each fixed $m>0$,
$q=m^2x^2/C^2$ is a local ordinary diffeomorphism there, so its
multiplicity in $q$ is four.

Apply the real scalar smooth preparation theorem
\cite[Theorem~5.3]{dencker1993}, treating $(m,\Theta)$ as parameters,
to $t=x-x_0$. It factors $F$ as a nonzero smooth function times
\[
 t^4+a_3(m,\Theta)t^3+a_2(m,\Theta)t^2
                         +a_1(m,\Theta)t+a_0(m,\Theta).
\]
The translation $t=y-a_3/4$ gives \eqref{quartic:preparation}.
At $m=0$, $\Theta=\Theta_0$, the factor is $-1$ and all
$a_j$ vanish. At $m=0$, write $\Delta\Theta=\Theta-\Theta_0$;
comparison of control derivatives modulo $t^3$ gives
\[
 \begin{aligned}
 u&=\Delta S+O(|\Delta\Theta|^2),\\
 v&=2x_0\Delta S-\Delta T+O(|\Delta\Theta|^2),\\
 w&=-\Delta A+x_0^2\Delta S-x_0\Delta T
                                    +O(|\Delta\Theta|^2).
 \end{aligned}
\]
The determinant of these three linear forms in $(\Delta A,\Delta S,
\Delta T)$ is $1$. Hence $(u,v,w)$ are actual local control
coordinates, jointly with $m$; equivalently,
$\det D_{(A,S,T)}(F,F_x,F_{xx})=2+O(m)\ne0$ on the branch.
Their origin corresponds to the actual quadruple point by the
uniqueness already proved. These changes identify zero equations
and their multiplicities; they do not assert conjugacy of return maps.

First fix this preparation neighbourhood and then choose the positive
interval in Lemma~\ref{mainloc:all-roots} inside its state chart.
Shrinking the common control neighbourhood makes every real root
of the monic quartic lie in that chart, since all its coefficients
vanish at the centre. Thus the prepared equation counts exactly
the whole positive cluster. At nondegenerate controls its real
simple-root count is $0$, $2$ or $4$. The polynomials
\[
 y^4+\eta^4,\qquad y^4-\eta^4,\qquad
 (y^2-\eta^2)(y^2-4\eta^2),\qquad 0<\eta\ll1,
\]
realize all three possibilities on open coefficient sets. The actual
control diffeomorphism transfers them to open sets in the family.
Lemma~\ref{mainloc:all-roots} supplies exactly two additional simple
cycles throughout. This gives the stated counts and, at the origin,
the three distinct cycles.
\end{proof}

\begin{proposition}[Local quartic strata]\label{quartic:strata}
In the actual preparation controls of \eqref{quartic:preparation},
the multiple-root set is parametrized by
\begin{equation}\label{quartic:double-surface}
 v=-4t^3-2ut,\qquad w=3t^4+ut^2.
\end{equation}
Its triple-root curve, its two-double-root curve and its quadruple
point are respectively
\[
 \begin{aligned}
 (u,v,w)&=(-6t^2,8t^3,-3t^4),&&t\ne0,\\
 (u,v,w)&=(-2t^2,0,t^4),&&t\ne0,\\
 (u,v,w)&=(0,0,0).
 \end{aligned}
\]
The ordinary fold part of \eqref{quartic:double-surface} excludes
both the triple-root curve and a simultaneous second double root.
These are exact local strata in the transformed actual controls.
\end{proposition}

\begin{proof}
Solving $p(t)=p'(t)=0$ gives \eqref{quartic:double-surface}.
On this surface, $p''(t)=0$ gives $p=(y-t)^3(y+3t)$,
and the additional equation $p'''(t)=0$ forces $t=0$.
Two distinct double roots must be $t,-t$, since the cubic coefficient
vanishes, so $p=(y^2-t^2)^2$.
\end{proof}

Figure~\ref{fig:quartic-local-diagram} displays the local root counts
in these preparation controls. The control change depends on $m$;
the strata above describe a neighbourhood of the actual quadruple
cycle, rather than the entire boundary of the main-scale domain.
% EXACT discriminant of the actual prepared quartic, on u=-eps^2.
% Similarity coordinates: y=eps Y, V=v/eps^3, W=w/eps^4.
% Curve parameter: (V,W)=(2t-4t^3,3t^4-t^2).
% Source and independent polynomial checks: scripts/check-local-diagrams.py.
\begin{figure}[!htbp]
\centering
\begin{tikzpicture}[x=6.6cm,y=4.7cm,font=\small,
  boundary/.style={black,line width=.75pt},
  guide/.style={black!55,line width=.35pt},
  leader/.style={black,line width=.4pt,-{Stealth[length=1.4mm]}},
  region/.style={fill=white,inner sep=1.5pt,font=\normalsize}]
  \begin{scope}
    \clip (-.94,-.15) rectangle (.94,.56);
    % The closed central chamber has four simple quartic roots.
    \fill[black!9]
      plot[domain=-.707106781186548:.707106781186548,
        variable=\t,samples=181]
      ({2*\t-4*(\t)^3},{3*(\t)^4-(\t)^2}) -- cycle;
    \draw[boundary]
      plot[domain=-.83:.83,variable=\t,samples=301]
      ({2*\t-4*(\t)^3},{3*(\t)^4-(\t)^2});
  \end{scope}
  \draw[guide,-{Stealth[length=1.5mm]}]
    (-.94,-.15)--(.99,-.15) node[right,black] {$V$};
  \draw[guide,-{Stealth[length=1.5mm]}]
    (-.94,-.15)--(-.94,.60) node[above,black] {$W$};
  \foreach \x/\lab in {-.5/{-\tfrac12},0/0,.5/{\tfrac12}}{
    \draw[guide] (\x,-.15)--(\x,-.164);
    \node[below,font=\scriptsize] at (\x,-.164) {$\lab$};
  }
  \foreach \y/\lab in {0/0,.25/{\tfrac14},.5/{\tfrac12}}{
    \draw[guide] (-.94,\y)--(-.955,\y);
    \node[left,font=\scriptsize] at (-.955,\y) {$\lab$};
  }
  \node[region] at (0,.46) {$2$};
  \node[region,fill=black!9] at (0,.12) {$6$};
  \node[region] at (-.76,.13) {$4$};
  \node[region] at (.76,.13) {$4$};
  \node[region] at (0,-.105) {$4$};
  % Two triple-root cusp points and the two-double-root crossing.
  \draw[boundary,fill=white]
    ({-4/(3*sqrt(6))},{-1/12}) circle[radius=1.6pt];
  \draw[boundary,fill=white]
    ({4/(3*sqrt(6))},{-1/12}) circle[radius=1.6pt];
  \fill (0,.25) circle[radius=1.8pt];
  \node[anchor=east,align=right,font=\scriptsize] (dd)
    at (-.29,.30) {two double\\roots};
  \draw[leader] (dd.east)--(-.015,.252);
  \node[anchor=west,font=\scriptsize] (fold) at (.50,.44) {fold};
  \draw[leader] (fold.west)--(.238,.424);
  \node[anchor=east,font=\scriptsize] (triple) at (-.66,-.055) {triple};
  \draw[leader] (triple.east)--(-.551,-.081);
\end{tikzpicture}
\caption{A slice of the actual prepared quartic
$y^4+uy^2+vy+w$ at $u=-\varepsilon^2$, for any sufficiently small
$\varepsilon>0$. In the similarity coordinates
$V=v/\varepsilon^3$, $W=w/\varepsilon^4$, the exact discriminant is
$(V,W)=(2t-4t^3,3t^4-t^2)$. Region labels are the total numbers of
simple cycles, including the two simple negative companions.
Open circles mark triple roots; the crossing has two double roots.
This slice does not contain the quadruple-root point $u=v=w=0$.
The boundary is exact in actual prepared controls.}
\label{fig:quartic-local-diagram}
\end{figure}


\FloatBarrier

\begin{corollary}[Open six-cycle regions]\label{cor:six-open}
Within every fixed permitted finite-dimensional family, the parameters
giving six simple fixed points in $I$ form an open subset of the
common parameter neighbourhood.
\end{corollary}

\begin{proof}
Fix a return with six simple zeros in $I$. The displacement is jointly
ordinary $C^1$ in state and parameters, including at its thresholds.
Simple-root continuation in six disjoint subintervals compactly
contained in $I$, together with the common six-point bound of
Proposition~\ref{prop:common-count}, proves openness.
\end{proof}
